\documentclass[11pt]{article}

\usepackage[T1]{fontenc}
\usepackage{lmodern}
\usepackage[margin=1in]{geometry}
\usepackage{amsmath}
\usepackage{amssymb}
\usepackage{booktabs}

\title{Text Property Graph Construction and the Multiview GNN Model}
\author{}
\date{}

\begin{document}
\maketitle

\section{Purpose}

Each CVE record is converted into a Text Property Graph (TPG), and the
multiview GNN consumes that graph to produce a prediction. This note
walks through the construction pipeline, the model architecture, and
the training settings. The goal is to be precise enough that an
engineer can reimplement the model, but light enough on notation that a
practitioner can read it in one sitting.

The model has two branches:

\begin{itemize}
    \item the \textbf{graph branch}, which learns from the TPG
    structure (nodes, typed edges, contextual node features);
    \item the \textbf{tabular branch}, which learns from CVE-level
    metadata (CVSS, CWE, age, source counts, public-exploit
    availability).
\end{itemize}

Edge types are graph-branch features. The tabular branch never sees
them.

\section{TPG Construction Procedure}

Each CVE graph is built from text. In the description-only baseline
the text is the CVE description. In the summary-only variants the text
is one of the LLM summary columns. In the combined (\textsc{ALL})
variant the text is the description plus all three summaries joined.

Construction proceeds in seven steps:

\begin{enumerate}
    \item Pick the text source (description, one summary, or
    description + summaries combined).
    \item Run the spaCy linguistic frontend. This produces document,
    paragraph, sentence, token, named-entity, predicate, argument,
    noun-phrase, verb-phrase and clause nodes, plus the base linguistic
    edges (\texttt{DEP}, \texttt{NEXT\_TOKEN}, \texttt{NEXT\_SENT},
    \texttt{NEXT\_PARA}, \texttt{COREF}, \texttt{SRL\_ARG},
    \texttt{CONTAINS}, \texttt{BELONGS\_TO}, \texttt{ENTITY\_REL},
    \texttt{RST\_RELATION}, \texttt{DISCOURSE}, \texttt{SIMILARITY},
    \texttt{AMR\_EDGE}).
    \item Run the security frontend. This adds CVE, CWE, software,
    version, attack-vector, impact, severity and remediation entities
    on top of the linguistic graph. With \texttt{--include-security-edges},
    the security-relations pass also emits ten typed \texttt{SEC\_*}
    edges between those entities.
    \item Compute a 768-dim contextual vector for every text-bearing
    node with SecBERT (jackaduma/SecBERT, frozen). SecBERT runs on
    WordPiece subtokens and the resulting subtoken vectors are
    averaged within the boundary of each node's text span (a token's
    own subtokens for token nodes; the constituent tokens for
    sentence, paragraph, entity, predicate, noun-phrase and clause
    nodes). Nodes without text get a zero vector. Every node ends up
    with the same 768-dim slot in its feature vector, regardless of
    type.
    \item Convert the graph to PyTorch Geometric tensors:
    \texttt{x} (node features), \texttt{edge\_index},
    \texttt{edge\_type}, and \texttt{edge\_attr} (one-hot edge type).
    \item Attach the target. In soft-label runs the target is the EPSS
    probability $\in [0, 1]$. In binary-mode runs it is either CISA
    KEV membership or the EPSS score thresholded at $0.1$.
\end{enumerate}

\section{Graph Representation}

After construction, a CVE graph has $N$ nodes, $M$ directed edges,
and an edge-type vocabulary of size $|R|$. The vocabulary is $|R| = 13$
for the base schema, or $|R| = 23$ when security edges are turned on
(the 13 base types plus the 10 \texttt{SEC\_*} types).

\paragraph{Node features.}
Each node $i$ has a feature vector

\[
x_i = \bigl[\,\text{node-type one-hot (13 dims)};\
              \text{SecBERT embedding (768 dims)}\,\bigr]
\;\in\; \mathbb{R}^{781}.
\]

If \texttt{--add-source-labels} is set, a 3-dim source one-hot
(description, summary, mixed) is appended, giving 784 dims. The 13
node types are document, paragraph, sentence, token, entity,
predicate, argument, concept, noun-phrase, verb-phrase, clause,
mention, topic.

\paragraph{Edge tensors.}
The $M$ edges are stored as two integer tensors:

\begin{itemize}
    \item \texttt{edge\_index}: a $2 \times M$ tensor of source and
    target node indices.
    \item \texttt{edge\_type}: an $M$-vector of relation indices in
    the range $0$ to $|R| - 1$. Indices $0$ to $12$ are the base
    types; indices $13$ to $22$ are the \texttt{SEC\_*} types when
    security edges are enabled.
\end{itemize}

\section{Multiview Edge Partition}

A single-view GNN treats every edge identically. The multiview
decomposition splits the edge vocabulary into disjoint groups called
views, each capturing a different kind of relation. The base
configuration uses four views; a fifth view is added when the
security edges are enabled.

\begin{table}[h]
\centering
\small
\begin{tabular}{ll}
\toprule
View & Edge types in the view \\
\midrule
Syntactic   & \texttt{DEP}, \texttt{CONTAINS}, \texttt{BELONGS\_TO} \\
Sequential  & \texttt{NEXT\_TOKEN}, \texttt{NEXT\_SENT}, \texttt{NEXT\_PARA} \\
Semantic    & \texttt{COREF}, \texttt{SRL\_ARG}, \texttt{AMR\_EDGE} \\
Discourse   & \texttt{RST\_RELATION}, \texttt{DISCOURSE}, \texttt{ENTITY\_REL}, \texttt{SIMILARITY} \\
Security    & all \texttt{SEC\_*} types (only present when $|R| = 23$) \\
\bottomrule
\end{tabular}
\caption{The four base views plus the optional security view, with
their edge types.}
\end{table}

For each view $v$, we keep only that view's edges. Call the resulting
edge subset $E_v$. When security edges are off, $E_{\text{security}}$
is empty and that view's encoder acts as the identity (no information
flows through it).

\section{Per-View Encoder}

Before any view-specific processing, the model lifts every node feature
into a hidden space of dimension $h = 128$ with a small input block: a
linear projection $W_{\text{in}}$, followed by LayerNorm, GELU and
Dropout. The projected node features are written $H^{(0)}$.

Each view has its own encoder, which is a Gated Graph Neural Network
(GGNN; Li et al., 2015). The encoder runs $K$ message-passing steps
restricted to that view's edges:

\[
H_v = \text{ViewEncoder}_v\bigl(H^{(0)},\; E_v\bigr)
\;\in\; \mathbb{R}^{N \times h}.
\]

A single GGNN step is a standard message-passing update: each node
sums linear messages from its in-neighbours, then a GRU cell combines
that sum with the node's previous hidden state. The same update is
unrolled $K$ times so a node's representation can pick up information
from neighbours up to $K$ hops away. We add a residual connection
plus a LayerNorm and Dropout around the encoder, and use $K = 3$ in
all reported runs.

\section{Node-Conditioned View Fusion}

Each node now has $V$ candidate embeddings, one per active view
($V = 4$ in the base configuration, $V = 5$ when the security view is
on). To combine them, we use a small attention mechanism that picks
how much weight each view gets, separately for each node.

For node $i$, the model computes a query vector $q_i$ from the input
projection and a key vector $k_i^{v}$ from each view's embedding.
Both vectors live in $\mathbb{R}^{h}$. The model scores each view by
the scaled dot product $q_i \cdot k_i^{v} / \sqrt{h}$, normalises the
$V$ scores across views with a softmax to get attention weights
$\alpha_i^{v}$, and forms the fused embedding as the weighted sum:

\[
h_i = \sum_{v}\, \alpha_i^{v}\, H_v[i, :].
\]

Because the attention weights are computed per node, different nodes
can rely on different views. A token inside a security entity can lean
on the security view; a discourse connective can lean on the discourse
view; a verb in a predicate-argument structure can lean on the
semantic view.

\section{Graph Pooling and Prediction Head}

The $N$ fused node embeddings are pooled into a single graph vector
$g$ by concatenating their element-wise mean and their element-wise
maximum across nodes. With $h = 128$, $g$ has $2h = 256$ dimensions.
A small MLP head then turns $g$ into a single logit $\hat{\ell}$, and
the sigmoid $\sigma(\hat{\ell})$ is the prediction. In the hybrid
configuration the head's input is the concatenation of $g$ with a
$64$-dim tabular embedding produced by a separate tabular MLP.

\section{Training Objective and Optimisation}

The loss is binary cross-entropy with logits between the predicted
probability $\sigma(\hat{\ell})$ and the target $y$. Positive examples
are up-weighted by a factor equal to the negative-to-positive ratio in
the training set, which keeps the rare high-EPSS class from being
drowned out. The same loss handles both hard $\{0, 1\}$ labels (KEV
membership) and soft $[0, 1]$ labels (raw EPSS scores), since BCE is
defined for any target in $[0, 1]$.

The optimiser is AdamW. Gradient norms are clipped at $1.0$ after
each backward pass. The learning rate is reduced on plateau (mode
max, factor $0.5$, patience $5$, minimum learning rate $10^{-6}$),
and training stops early after $15$ epochs without an improvement in
validation PR-AUC. The threshold $0.1$ on EPSS is used only for
stratified splitting and for binarising the labels at metric time
(precision, recall and F1 at decision threshold $0.5$). It is never
the training target.

\section{Summary of Hyperparameters}

The values used in the 20-run text-source ablation and the NVD/KEV
runs are collected here so they live next to the architecture
description.

\begin{table}[h]
\centering
\small
\begin{tabular}{lll}
\toprule
Component & Setting & Value \\
\midrule
Backbone                   & \texttt{--backbone}        & multiview \\
Hidden dimension           & $h$                        & 128 \\
Tabular MLP width          & $h_t$                      & 64 \\
GGNN message-passing steps & $K$ (\texttt{--layers})    & 3 \\
Dropout                    & \texttt{--dropout}         & 0.3 \\
Epochs                     & \texttt{--epochs}          & 100 \\
Batch size                 & \texttt{--batch-size}      & 32 \\
Learning rate              & \texttt{--lr}              & $1 \times 10^{-3}$ \\
Weight decay               & \texttt{--weight-decay}    & $1 \times 10^{-4}$ \\
Gradient-norm clip         &                            & 1.0 \\
Scheduler                  & reduce-on-plateau          & mode max, factor 0.5, patience 5, min lr $10^{-6}$ \\
Early-stopping patience    & \texttt{--patience}        & 15 epochs without val PR-AUC gain \\
Random seed                & \texttt{--seed}            & 42 \\
SecBERT embedding dim      & \texttt{--embedding-dim}   & 768 \\
Top-K CWEs (NVD/KEV)       & \texttt{--top-k-cwes}      & 25 \\
\bottomrule
\end{tabular}
\caption{Hyperparameters used in the reported experiments.}
\end{table}

\section{What the Architecture Actually Encodes}

Edge labels are not just counted. Their placement in the graph is part
of the model input. A \texttt{SEC\_AFFECTS} edge tells the GNN which
CVE node connects to which software node. A tabular one-hot feature
could only record that such a relation exists somewhere, losing the
structural information.

For this reason \texttt{edge\_index} together with \texttt{edge\_type}
participate in typed message passing through the per-view GGNN
encoders, while tabular features carry only document-level metadata.

End-to-end, the graph branch turns $X$, \texttt{edge\_index} and
\texttt{edge\_type} into a 256-dim graph vector $g$ via input
projection, $V$ per-view GGNN encoders, node-conditioned view fusion,
and mean-plus-max pooling. The tabular branch turns the tabular
feature vector into a 64-dim embedding via a small MLP. In the hybrid
configuration the prediction is the sigmoid of a fusion MLP applied
to the concatenation of the graph vector and the tabular embedding;
in the graph-only configuration the prediction comes from $g$ alone.

\end{document}
